\section{Three powers and one common prime}\label{sec:statement}
\begin{conjecture}[Beal]\label{conj:beal}
If positive integers $A,B,C$ and integers $x,y,z>2$ satisfy
\begin{equation}\label{eq:beal}
A^x+B^y=C^z,
\end{equation}
then $A,B,C$ have a common prime divisor.
\end{conjecture}
The exponents are allowed to differ. That is the essential extra freedom beyond the equal exponent Fermat equation. The common prime belongs to all three bases, not merely to the product of two of them. The historical statement and its generalized Fermat context are discussed in Mauldin's account and its authenticated published translation \cite{mauldin1997,mauldin1998translation}.

\begin{definition}
A positive solution of \eqref{eq:beal} is \emph{primitive} when $\gcd(A,B,C)=1$. Its ordered exponent triple $(x,y,z)$ is its signature. The first two entries may be exchanged together with the first two bases.
\end{definition}
Primitivity becomes stronger on an actual solution than it is on an arbitrary integer triple.
\begin{lemma}\label{lem:pairwise}
For a solution of \eqref{eq:beal}, every prime dividing two bases divides the third. Consequently a primitive solution has pairwise coprime bases.
\end{lemma}
\begin{proof}
If a prime $p$ divides $A$ and $B$, the equation gives $p\mid C^z$, hence $p\mid C$. If $p$ divides $A$ and $C$, it gives $p\mid B^y$, hence $p\mid B$. The third case is the same. Thus a prime shared by any pair would divide the gcd of all three.
\end{proof}
Beal is therefore equivalent to the absence of primitive positive solutions for all signatures with entries greater than two. This reformulation is the one used throughout the research program. It makes the support of each base disjoint from the other two supports, which is useful when assigning local valuation sectors and conductor conditions.

\begin{example}\label{ex:commonprime}
The equality $3^3+6^3=3^5$ is a valid power equation. Both terms on the left are positive, and all exponents exceed two. The common prime is $3$, so this equality illustrates the conjecture rather than contradicting it. Its powered terms are $27,216,243$.
\end{example}
\figplaceholder{01_beal_terms}{The equation in Example~\ref{ex:commonprime} has a common prime in its bases. Exact equality alone does not supply primitivity.}
An ordinary division by $\gcd(A,B,C)$ is not a valid reduction to a primitive solution. If $A=dA'$, $B=dB'$, $C=dC'$, the three terms carry factors $d^x,d^y,d^z$. Unless the exponents agree, division by one of these powers changes the other coefficients. Later, a weighted scaling will preserve the equation. That operation has a precise valuation criterion and is not the same as dividing every base by $d$.

The positivity assumptions also carry information. They imply
\begin{equation}\label{eq:heightbasic}
A^x<C^z,\qquad B^y<C^z,\qquad C\ge2.
\end{equation}
The last inequality follows because $A^x+B^y\ge2$. A base on the left may equal $1$, so a finite classification of powered values need not be a finite classification of exponent labels. The equality $1^x=1$ holds for every $x$. This small point is relevant when comparing a theorem for one fixed signature with a statement uniform in all exponents.

No bounded search can settle Conjecture~\ref{conj:beal} without a theorem covering its unsearched range. A search can exclude a stated finite region, identify local obstructions, and check an implementation. A reduction to finitely many residual patterns can be stronger than an ordinary integer search, but it still needs an exhaustive arithmetic cover and a proof that every covered pattern is impossible.
